\section{Evaluation}
\label{sec:eval}
% We conduct empirical evaluations to validate the effectiveness of the proposed \model\ framework. The experiments aim to evaluate \model\ from the following dimensions:
% $\bullet$ How does \model\ perform in comparison to existing RAG baseline methods?
% $\bullet$ How does the dual-level retrieval and the graph-based indexing benefit the generation quality of \model?
% $\bullet$ What specific advantages does \model\ exhibit in specific cases?
% $\bullet$ How is the cost and the adaptation abilities of \model?
We conduct empirical evaluations on benchmark data to assess the effectiveness of the proposed \model\ framework by addressing the following research questions: $\bullet$ \textbf{(RQ1)}: How does \model\ compare to existing RAG baseline methods in terms of generation performance? $\bullet$ \textbf{(RQ2)}: How do dual-level retrieval and graph-based indexing enhance the generation quality of \model? $\bullet$ \textbf{(RQ3)}: What specific advantages does \model\ demonstrate through case examples in various scenarios? $\bullet$ \textbf{(RQ4)}: What are the costs associated with \model, as well as its adaptability to data changes?

\subsection{Experimental Settings}
\textbf{Evaluation Datasets}.
% To conduct a comprehensive analysis of \model, we selected four datasets from the UltraDomain benchmark~\citep{qian2024memorag}. The UltraDomain data is sourced from 428 college textbooks, covering 18 distinct domains, including agriculture, social sciences, and humanities. From these domains, we chose the Agriculture, CS, Legal, and Mix datasets. The total number of tokens in each dataset ranges from 600,000 to 5,000,000, with detailed information provided in Table~\ref{dataset}. Below is a specific introduction to the four domains used in our experiments:
To conduct a comprehensive analysis of \model, we selected four datasets from the UltraDomain benchmark~\citep{qian2024memorag}. The UltraDomain data is sourced from 428 college textbooks and encompasses 18 distinct domains, including agriculture, social sciences, and humanities. From these, we chose the Agriculture, CS, Legal, and Mix datasets. Each dataset contains between 600,000 and 5,000,000 tokens, with detailed information provided in Table~\ref{dataset}. Below is a specific introduction to the four domains utilized in our experiments:
\begin{itemize}[leftmargin=*]
    \item \textbf{Agriculture}: This domain focuses on agricultural practices, covering a range of topics including beekeeping, hive management, crop production, and disease prevention.
    
    \item  \textbf{CS}: This domain focuses on computer science and encompasses key areas of data science and software engineering. It particularly highlights machine learning and big data processing, featuring content on recommendation systems, classification algorithms, and real-time analytics using Spark.
    
    \item \textbf{Legal}: This domain centers on corporate legal practices, addressing corporate restructuring, legal agreements, regulatory compliance, and governance, with a focus on the legal and financial sectors.
    
    \item \textbf{Mixed}: This domain presents a rich variety of literary, biographical, and philosophical texts, spanning a broad spectrum of disciplines, including cultural, historical, and philosophical studies.
\end{itemize}

\textbf{Question Generation}. 
% To evaluate the effectiveness of RAG systems for more high-level sensemaking tasks, we merge all the text content from each dataset as context and follow the generation method used in~\cite{edge2024local}. Specifically, we task an LLM with generating 5 RAG users and 5 tasks for each user. Each generated user is associated with a textual description for his/her expertise and traits that motivates his/her activities in raising questions. Each task of users is also described with texts, highlighting one of the user's potential intentions when interacting with RAG systems. For each user-task combination, the LLM is then required to generate 5 questions that necessitate an understanding of the entire corpus. In total, 125 questions are generated for each dataset.
To evaluate the effectiveness of RAG systems for high-level sensemaking tasks, we consolidate all text content from each dataset as context and adopt the generation method outlined in~\cite{edge2024local}. Specifically, we instruct an LLM to generate five RAG users, along with five tasks for each user. Each generated user is accompanied by a textual description detailing their expertise and traits that motivate their question-raising activities. Each user task is also described, emphasizing one of the user's potential intentions when interacting with RAG systems. For each user-task combination, the LLM generates five questions that require an understanding of the entire corpus. In total, this process results in 125 questions for each dataset.

\textbf{Baselines}.  \model\ is compared against the following state-of-the-art methods across all datasets:
\begin{itemize}[leftmargin=*]
    \item \textbf{Naive RAG}~\citep{gao2023retrieval}: This model serves as a standard baseline in existing RAG systems. It segments raw texts into chunks and stores them in a vector database using text embeddings. For queries, Naive RAG generates vectorized representations to directly retrieve text chunks based on the highest similarity in their representations, ensuring efficient and straightforward matching.
    % This model is a standard baseline in simple RAG systems. It segments raw texts into chunks and stores them in a vector database with text embeddings. For queries, Naive RAG creates vectorized representations to directly retrieve text chunks based on the highest similarity in their vector representations, ensuring efficient and straightforward matching.
    
    \item \textbf{RQ-RAG}~\citep{chan2024rq}: This approach leverages the LLM to decompose the input query into multiple sub-queries. These sub-queries are designed to enhance search accuracy by utilizing explicit techniques such as rewriting, decomposition, and disambiguation.
    % This approach utilizes the LLM to decompose the input query into several sub-queries. These sub-queries are designed to enhance search accuracy through explicit query rewriting, decomposition, and disambiguation.
    
    \item \textbf{HyDE}~\citep{hyde}: This method utilizes the LLM to generate a hypothetical document based on the input query. This generated document is then employed to retrieve relevant text chunks, which are subsequently used to formulate the final answer.
    
    \item \textbf{GraphRAG}~\citep{edge2024local}: This is a graph-enhanced RAG system that utilizes an LLM to extract entities and relationships from the text, representing them as nodes and edges. It generates corresponding descriptions for these elements, aggregates nodes into communities, and produces a community report to capture global information. When handling high-level queries, GraphRAG retrieves more comprehensive information by traversing these communities.
\end{itemize}

\textbf{Implementation and Evaluation Details}.
In our experiments, we utilize the \href{https://github.com/gusye1234/nano-vectordb}{\textit{nano vector database}} for vector data management and access. For all LLM-based operations in \model, we default to using GPT-4o-mini. To ensure consistency, the chunk size is set to 1200 across all datasets. Additionally, the gleaning parameter is fixed at 1 for both GraphRAG and \model.

% As it is hard to define ground truth for many RAG queries, especially for complex queries involving high-level semantics, we follow existing work~\citep{edge2024local} and adopt the LLM-based multi-dimensional comparison method. A strong LLM, specifically GPT-4o-mini, is used to rank between each baseline and our \model. The prompt we used for evaluation is elaborated in Appendix~\ref{app:evaluation_prompt}. Totally four evaluation dimensions are utilized, including:
Defining ground truth for many RAG queries, particularly those involving complex high-level semantics, poses significant challenges. To address this, we build on existing work~\citep{edge2024local} and adopt an LLM-based multi-dimensional comparison method. We employ a robust LLM, specifically GPT-4o-mini, to rank each baseline against our \model. The evaluation prompt we used is detailed in Appendix~\ref{app:evaluation_prompt}. In total, we utilize four evaluation dimensions, including:

% i) \textbf{Comprehensiveness}: How much detail does the answer provide to cover all aspects and details of the question?
% ii) \textbf{Diversity}: How varied and rich is the answer in providing different perspectives and insights on the question?
% iii) \textbf{Empowerment}: How well does the answer help the reader understand and make informed judgments about the topic?
% % iv) \textbf{Directness}: How specifically does the answer address the question?
% iv) \textbf{Overall}: Considers the cumulative performance across the other three dimensions to determine the best overall answer.
i) \textbf{Comprehensiveness}: How thoroughly does the answer address all aspects and details of the question? ii) \textbf{Diversity}: How varied and rich is the answer in offering different perspectives and insights related to the question? iii) \textbf{Empowerment}: How effectively does the answer enable the reader to understand the topic and make informed judgments? iv) \textbf{Overall}: This dimension assesses the cumulative performance across the three preceding criteria to identify the best overall answer.

% The LLM directly compare two answers for each dimension and select the superior answer for each dimension. After determining the winning answer for each of the three dimensions, the LLM then combine the results across all dimensions to identify the overall better answer. To avoid any bias caused by the order of answers in the prompt during evaluation, we alternate the placement of each answer in the prompt and calculate win rates accordingly to obtain the final result.
The LLM directly compares two answers for each dimension and selects the superior response for each criterion. After identifying the winning answer for the three dimensions, the LLM combines the results to determine the overall better answer. To ensure a fair evaluation and mitigate the potential bias that could arise from the order in which the answers are presented in the prompt, we alternate the placement of each answer. We calculate win rates accordingly, ultimately leading to the final results. 

\subsection{Comparison of \model\ with Existing RAG Methods (RQ1)}

% We compare \model\ against each baseline one-by-one on the four evaluation dimensions and four datasets. The results are reported in Table~\ref{tab:performance}. From the results, we make the following conclusions:
We compare \model\ against each baseline across various evaluation dimensions and datasets. The results are presented in Table~\ref{tab:performance}. Based on these findings, we draw the following conclusions:

\begin{table}[t]
\centering
\caption{Win rates (\%) of baselines v.s. \model\ across four datasets and four evaluation dimensions.}
\label{tab:performance}
\vspace{-0.1in}
% \setlength{\tabcolsep}{1.8mm}
% \scriptsize
\resizebox{\textwidth}{!}{
\begin{tabular}{@{}lcccccccc@{}}
\toprule
\textbf{}    & \multicolumn{2}{c}{\textbf{Agriculture}} & \multicolumn{2}{c}{\textbf{CS}} & \multicolumn{2}{c}{\textbf{Legal}} & \multicolumn{2}{c}{\textbf{Mix}} \\ 
\cmidrule(lr){2-3} \cmidrule(lr){4-5} \cmidrule(lr){6-7} \cmidrule(lr){8-9}
                      & NaiveRAG & \textbf{LightRAG} & NaiveRAG & \textbf{LightRAG} & NaiveRAG & \textbf{LightRAG} & NaiveRAG & \textbf{LightRAG} \\
\midrule
Comprehensiveness      & 32.4\%      & \underline{67.6\%}     & 38.4\%      & \underline{61.6\%}     & 16.4\%      & \underline{83.6\%}     & 38.8\%      & \underline{61.2\%}     \\
Diversity              & 23.6\%      & \underline{76.4\%}     & 38.0\%      & \underline{62.0\%}     & 13.6\%      & \underline{86.4\%}     & 32.4\%      & \underline{67.6\%}     \\
Empowerment            & 32.4\%      & \underline{67.6\%}     & 38.8\%      & \underline{61.2\%}     & 16.4\%      & \underline{83.6\%}     & 42.8\%      & \underline{57.2\%}     \\
% Directness             & \underline{70.98\%}      & 29.02\%     & \underline{59.84\%}      & 40.16\%     & \underline{55.69\%}      & 44.31\%     & \underline{69.95\%}      & 30.05\%     \\
Overall                & 32.4\%      & \underline{67.6\%}     & 38.8\%      & \underline{61.2\%}     & 15.2\%      & \underline{84.8\%}     & 40.0\%      & \underline{60.0\%}     \\
\cmidrule(lr){2-3} \cmidrule(lr){4-5} \cmidrule(lr){6-7} \cmidrule(lr){8-9}
                      & RQ-RAG & \textbf{LightRAG} & RQ-RAG & \textbf{LightRAG} & RQ-RAG & \textbf{LightRAG} & RQ-RAG & \textbf{LightRAG} \\
\midrule
Comprehensiveness      & 31.6\%      & \underline{68.4\%}     & 38.8\%      & \underline{61.2\%}     & 15.2\%      & \underline{84.8\%}     & 39.2\%      & \underline{60.8\%}     \\
Diversity              & 29.2\%      & \underline{70.8\%}     & 39.2\%      & \underline{60.8\%}     & 11.6\%      & \underline{88.4\%}     & 30.8\%      & \underline{69.2\%}     \\
Empowerment            & 31.6\%      & \underline{68.4\%}     & 36.4\%      & \underline{63.6\%}     & 15.2\%      & \underline{84.8\%}     & 42.4\%      & \underline{57.6\%}     \\
% Directness             & \underline{60.53\%}      & 39.47\%     & \underline{62.32\%}      & 37.68\%     & \underline{51.15\%}      & 48.85\%     & \underline{71.14\%}      & 28.86\%     \\
Overall         & 32.4\%      & \underline{67.6\%}     & 38.0\%      & \underline{62.0\%}     & 14.4\%      & \underline{85.6\%}     & 40.0\%      & \underline{60.0\%}     \\
\cmidrule(lr){2-3} \cmidrule(lr){4-5} \cmidrule(lr){6-7} \cmidrule(lr){8-9}
                      & HyDE & \textbf{LightRAG} & HyDE & \textbf{LightRAG} & HyDE & \textbf{LightRAG} & HyDE & \textbf{LightRAG} \\
\midrule
Comprehensiveness      & 26.0\%      & \underline{74.0\%}     & 41.6\%      & \underline{58.4\%}     & 26.8\%      & \underline{73.2\%}     & 40.4\%      & \underline{59.6\%}     \\
Diversity              & 24.0\%      & \underline{76.0\%}     & 38.8\%      & \underline{61.2\%}     & 20.0\%      & \underline{80.0\%}     & 32.4\%      & \underline{67.6\%}     \\
Empowerment            & 25.2\%      & \underline{74.8\%}     & 40.8\%      & \underline{59.2\%}     & 26.0\%      & \underline{74.0\%}     & 46.0\%      & \underline{54.0\%}     \\
% Directness             & \underline{62.86\%}      & 37.14\%     & \underline{67.69\%}      & 32.31\%     & \underline{63.69\%}      & 36.31\%     & \underline{69.89\%}      & 30.11\%     \\
Overall         & 24.8\%      & \underline{75.2\%}     & 41.6\%      & \underline{58.4\%}     & 26.4\%      & \underline{73.6\%}     & 42.4\%      & \underline{57.6\%}     \\
\cmidrule(lr){2-3} \cmidrule(lr){4-5} \cmidrule(lr){6-7} \cmidrule(lr){8-9}
                      & GraphRAG & \textbf{LightRAG} & GraphRAG & \textbf{LightRAG} & GraphRAG & \textbf{LightRAG} & GraphRAG & \textbf{LightRAG} \\
\midrule
Comprehensiveness      & 45.6\%      & \underline{54.4\%}     & 48.4\%      & \underline{51.6\%}     &  48.4\%     & \underline{51.6\%}     & \underline{50.4\%}      & 49.6\%     \\
Diversity              & 22.8\%      & \underline{77.2\%}     & 40.8\%      & \underline{59.2\%}     &  26.4\%     & \underline{73.6\%}     & 36.0\%      & \underline{64.0\%}     \\
Empowerment            & 41.2\%      & \underline{58.8\%}     & 45.2\%      & \underline{54.8\%}     & 43.6\%      & \underline{56.4\%}     & \underline{50.8\%}      & 49.2\%     \\
% Directness             & \underline{69.85\%}      & 30.15\%     & \underline{56.96\%}      & 43.04\%     & \underline{58.11\%}      & 41.89\%     & \underline{55.44\%}      & 44.56\%     \\
Overall                & 45.2\%      & \underline{54.8\%}     & 48.0\%      & \underline{52.0\%}     & 47.2\%      & \underline{52.8\%}     & \underline{50.4\%}      & 49.6\%     \\
\bottomrule
\end{tabular}
}
\end{table}

\textbf{The Superiority of Graph-enhanced RAG Systems in Large-Scale Corpora}
% \textbf{Advantage of Graph-enhanced Retrieval in Large-Scale Token Queries}: 
% When dealing with large token counts and complex queries requiring comprehensive consideration of the dataset's background, Graph-based RAG systems like \model\ and GraphRAG consistently outperform chunk-based retrieval methods (NaiveRAG, HyDE, and RQRAG). This performance gap becomes particularly evident as the dataset size increases. For instance, in the smallest dataset (Mix), the performance difference between chunk-based methods and \model\ is relatively modest. However, in the largest dataset (Legal), the gap widens significantly, with chunk-based methods only achieving around 20\% \% win rates overall compared to \model's dominance. This trend highlights the scalability advantage of Graph-based retrieval, especially in handling more extensive and intricate datasets.
When handling large token counts and complex queries that require a thorough understanding of the dataset's context, graph-based RAG systems like \model\ and GraphRAG consistently outperform purely chunk-based retrieval methods such as NaiveRAG, HyDE, and RQRAG. This performance gap becomes particularly pronounced as the dataset size increases. For instance, in the largest dataset (Legal), the disparity widens significantly, with baseline methods achieving only about 20\% win rates compared to the dominance of \model. This trend underscores the advantages of graph-enhanced RAG systems in capturing complex semantic dependencies within large-scale corpora, facilitating a more comprehensive understanding of knowledge and leading to improved generalization performance.
% This trend underscores the scalability advantage of graph-based retrieval, particularly in managing larger and more intricate datasets.

\textbf{Enhancing Response Diversity with \model}:
% Compared to all kinds of baselines, \model\ shows a noticeable edge in the Diversity metric, especially in the larger Legal dataset. Its consistent lead in Diversity suggests that it is more effective at generating a broader range of responses, particularly in scenarios where diverse content is critical. We attribute this advantage to the comprehensive information retrieval capabilities of \model\ based on the effective graph-based text indexing, which always successfully fetches the full context for queries.
Compared to various baselines, \model\ demonstrates a significant advantage in the Diversity metric, particularly within the larger Legal dataset. Its consistent lead in this area underscores \model's effectiveness in generating a wider range of responses, especially in scenarios where diverse content is essential. We attribute this advantage to \model's dual-level retrieval paradigm, which facilitates comprehensive information retrieval from both low-level and high-level dimensions. This approach effectively leverages graph-based text indexing to consistently capture the full context in response to queries.

\textbf{\model's Superiority over GraphRAG}: 
% Although both \model\ and GraphRAG leverage graph-based retrieval mechanisms, \model\ consistently outperforms GraphRAG across multiple dimensions, particularly in larger and more complex datasets. In the Agriculture, CS, and Legal datasets, each containing millions of tokens, \model\ achieved a clear and consistent lead, comprehensively surpassing GraphRAG. This demonstrates \model's strong global control over the datasets, reinforcing its advantage in diverse and complex data environments. By generating keywords during both graph creation and query processing, \model\ ensures that the graph remains extensible, efficiently handling a larger scope of data while maintaining performance.
While both \model\ and GraphRAG use graph-based retrieval mechanisms, \model\ consistently outperforms GraphRAG, particularly in larger datasets with complex language contexts. In the Agriculture, CS, and Legal datasets—each containing millions of tokens—\model\ shows a clear advantage, significantly surpassing GraphRAG and highlighting its strength in comprehensive information understanding within diverse environments. \textbf{Enhanced Response Variety}: By integrating low-level retrieval of specific entities with high-level retrieval of broader topics, \model\ boosts response diversity. This dual-level mechanism effectively addresses both detailed and abstract queries, ensuring a thorough grasp of information. \textbf{Complex Query Handling}: This approach is especially valuable in scenarios requiring diverse perspectives. By accessing both specific details and overarching themes, \model\ adeptly responds to complex queries involving interconnected topics, providing contextually relevant answers.

\subsection{Ablation Studies (RQ2)}
\begin{table}[t]
\centering
\caption{Performance of ablated versions of \model, using NaiveRAG as reference.}
\label{tab:ablation}
\vspace{-0.1in}
\setlength{\tabcolsep}{1.2\tabcolsep}
\resizebox{\textwidth}{!}{
\begin{tabular}{@{}lcccccccc@{}}
\toprule
\textbf{}    & \multicolumn{2}{c}{\textbf{Agriculture}} & \multicolumn{2}{c}{\textbf{CS}} & \multicolumn{2}{c}{\textbf{Legal}} & \multicolumn{2}{c}{\textbf{Mix}} \\ 
\cmidrule(lr){2-3} \cmidrule(lr){4-5} \cmidrule(lr){6-7} \cmidrule(lr){8-9}
                      & NaiveRAG & \textbf{\model} & NaiveRAG & \textbf{\model} & NaiveRAG & \textbf{\model} & NaiveRAG & \textbf{\model} \\
\midrule
Comprehensiveness      & 32.4\%      & \underline{67.6\%}     & 38.4\%      & \underline{61.6\%}     & 16.4\%      & \underline{83.6\%}     & 38.8\%      & \underline{61.2\%}     \\
Diversity              & 23.6\%      & \underline{76.4\%}     & 38.0\%      & \underline{62.0\%}     & 13.6\%      & \underline{86.4\%}     & 32.4\%      & \underline{67.6\%}     \\
Empowerment            & 32.4\%      & \underline{67.6\%}     & 38.8\%      & \underline{61.2\%}     & 16.4\%      & \underline{83.6\%}     & 42.8\%      & \underline{57.2\%}     \\
% Directness             & \underline{70.98\%}      & 29.02\%     & \underline{59.84\%}      & 40.16\%     & \underline{55.69\%}      & 44.31\%     & \underline{69.95\%}      & 30.05\%     \\
Overall                & 32.4\%      & \underline{67.6\%}     & 38.8\%      & \underline{61.2\%}     & 15.2\%      & \underline{84.8\%}     & 40.0\%      & \underline{60.0\%}     \\
\cmidrule(lr){2-3} \cmidrule(lr){4-5} \cmidrule(lr){6-7} \cmidrule(lr){8-9}
                      & NaiveRAG & \textbf{-High} & NaiveRAG & \textbf{-High} & NaiveRAG & \textbf{-High} & NaiveRAG & \textbf{-High} \\
\midrule
Comprehensiveness      & 34.8\%      & \underline{65.2\%}     & 42.8\%      & \underline{57.2\%}     & 23.6\%      & \underline{76.4\%}     & 40.4\%      & \underline{59.6\%}     \\
Diversity              & 27.2\%      & \underline{72.8\%}     & 36.8\%      & \underline{63.2\%}     & 16.8\%      & \underline{83.2\%}     & 36.0\%      & \underline{64.0\%}     \\
Empowerment            & 36.0\%      & \underline{64.0\%}     & 42.4\%      & \underline{57.6\%}     & 22.8\%      & \underline{77.2\%}     & 47.6\%      & \underline{52.4\%}     \\
% Directness             & \underline{67.90\%}      & 32.10\%     & \underline{70.18\%}      & 29.82\%     & \underline{51.57\%}      & 48.43\%     & \underline{63.96\%}      & 36.04\%     \\
Overall                & 35.2\%      & \underline{64.8\%}     & 44.0\%      & \underline{56.0\%}     & 22.0\%      & \underline{78.0\%}     & 42.4\%      & \underline{57.6\%}     \\
\cmidrule(lr){2-3} \cmidrule(lr){4-5} \cmidrule(lr){6-7} \cmidrule(lr){8-9}
                      & NaiveRAG & \textbf{-Low} & NaiveRAG & \textbf{-Low} & NaiveRAG & \textbf{-Low} & NaiveRAG & \textbf{-Low} \\
\midrule
Comprehensiveness      & 36.0\%      & \underline{64.0\%}     & 43.2\%      & \underline{56.8\%}     & 19.2\%      & \underline{80.8\%}     & 36.0\%      & \underline{64.0\%}     \\
Diversity              & 28.0\%      & \underline{72.0\%}     & 39.6\%      & \underline{60.4\%}     & 13.6\%      & \underline{86.4\%}     & 33.2\%      & \underline{66.8\%}     \\
Empowerment            & 34.8\%      & \underline{65.2\%}     & 42.8\%      & \underline{57.2\%}     & 16.4\%      & \underline{83.6\%}     & 35.2\%      & \underline{64.8\%}     \\
% Directness             & \underline{71.19\%}      & 28.81\%     & \underline{60.17\%}      & 39.83\%     & \underline{57.17\%}      & 42.83\%     & \underline{69.41\%}      & 30.59\%     \\
Overall                & 34.8\%      & \underline{65.2\%}     & 43.6\%      & \underline{56.4\%}     & 18.8\%      & \underline{81.2\%}     & 35.2\%      & \underline{64.8\%}     \\

\cmidrule(lr){2-3} \cmidrule(lr){4-5} \cmidrule(lr){6-7} \cmidrule(lr){8-9}
                      & NaiveRAG & \textbf{-Origin} & NaiveRAG & \textbf{-Origin} & NaiveRAG & \textbf{-Origin} & NaiveRAG & \textbf{-Origin} \\
\midrule
Comprehensiveness      & 24.8\%      & \underline{75.2\%}     & 39.2\%      & \underline{60.8\%}     & 16.4\%      & \underline{83.6\%}     & 44.4\%      & \underline{55.6\%}     \\
Diversity              & 26.4\%      & \underline{73.6\%}     & 44.8\%      & \underline{55.2\%}     & 14.4\%      & \underline{85.6\%}     & 25.6\%      & \underline{74.4\%}     \\
Empowerment            & 32.0\%      & \underline{68.0\%}     & 43.2\%      & \underline{56.8\%}     & 17.2\%      & \underline{82.8\%}     & 45.2\%      & \underline{54.8\%}     \\
% Directness             & \underline{61.94\%}      & 38.06\%     & \underline{62.07\%}      & 37.93\%     & \underline{54.51\%}      & 45.49\%     & \underline{63.33\%}      & 36.67\%     \\
Overall                & 25.6\%      & \underline{74.4\%}     & 39.2\%      & \underline{60.8\%}     & 15.6\%      & \underline{84.4\%}     & 44.4\%      & \underline{55.6\%}     \\
\bottomrule
\end{tabular}
}
\end{table}

% We further conduct ablation studies to validate the impact of our dual-level retrieval paradigm and the involvement of original text in our \model. The results are presented in Table~\ref{tab:ablation}.
We also conduct ablation studies to evaluate the impact of our dual-level retrieval paradigm and the effectiveness of our graph-based text indexing in \model. The results are presented in Table~\ref{tab:ablation}.

\textbf{Effectiveness of Dual-level Retrieval Paradigm}.
% We first examine the impact of low-level and high-level retrieval. Two ablated models, each lacking one of the modules, are compared with \model\ across four datasets. We make the following observations for different variants.
We begin by analyzing the effects of low-level and high-level retrieval paradigms. We compare two ablated models—each omitting one module—against \model\ across four datasets. Here are our key observations for the different variants:

\begin{itemize}[leftmargin=*]
    % \item \textbf{Low-level-only Retrieval}: The -High variant removes the high-order retrieval. This causes a significant performance drop for almost all datasets and metrics. This is mainly due to its low-level-only retrieval mechanism, which excessively focuses on entities and their direct neighbors. While allowing for deeper exploration of directly-related entities, it struggles to collect information for complex and difficult queries that require comprehensive insights.
    \item \textbf{Low-level-only Retrieval}: The -High variant removes high-order retrieval, leading to a significant performance decline across nearly all datasets and metrics. This drop is mainly due to its emphasis on the specific information, which focuses excessively on entities and their immediate neighbors. While this approach enables deeper exploration of directly related entities, it struggles to gather information for complex queries that demand comprehensive insights.

    % \item \textbf{High-level-only Retrieval}: The -Low variant focuses on capturing a broader range of content by leveraging entity-wise relationships instead of specific entities. This approach provides a clear advantage in the comprehensiveness dimension, as it is able to gather more extensive and varied information. However, the trade-off is a lack of depth in examining specific entities, which can limit its ability to provide highly detailed insights. As a result, this high-level-only retrieval method may struggle with tasks that demand more precise, detailed answers.
    \item \textbf{High-level-only Retrieval}: The -Low variant prioritizes capturing a broader range of content by leveraging entity-wise relationships rather than focusing on specific entities. This approach offers a significant advantage in comprehensiveness, allowing it to gather more extensive and varied information. However, the trade-off is a reduced depth in examining specific entities, which can limit its ability to provide highly detailed insights. Consequently, this high-level-only retrieval method may struggle with tasks that require precise, detailed answers.

    % \item \textbf{Hybird Mode}: The hybird mode, \textit{i.e.} the full version of \model, combines the strengths of both low-level and high-level modes, retrieving a broader set of relationships while simultaneously conducting an in-depth exploration of specific entities. This dual-level approach ensures both breadth in the retrieval process and depth in the analysis, offering a comprehensive view of the data. As a result, \model\ achieves a balanced performance across multiple dimensions.
    \item \textbf{Hybrid Mode}: The hybrid mode, or the full version of \model, combines the strengths of both low-level and high-level retrieval methods. It retrieves a broader set of relationships while simultaneously conducting an in-depth exploration of specific entities. This dual-level approach ensures both breadth in the retrieval process and depth in the analysis, providing a comprehensive view of the data. As a result, \model\ achieves balanced performance across multiple dimensions.
    
\end{itemize}

\textbf{Semantic Graph Excels in RAG}.
We eliminated the use of original text in our retrieval process. Surprisingly, the resulting variant, -Origin, does not exhibit significant performance declines across all four datasets. In some cases, this variant even shows improvements (\textit{e.g.} in Agriculture and Mix). We attribute this phenomenon to the effective extraction of key information during the graph-based indexing process, which provides sufficient context for answering queries. Additionally, the original text often contains irrelevant information that can introduce noise in the response.
% We eliminated the use of original text in our retrieval process. Surprisingly, the resulting variant, -Origin, does not show significant performance declines across all four datasets. In some instances, this variant even shows performance improvements (e.g., in Agriculture and Mix). We attribute this phenomenon to the effective extraction of key information during the graph-based indexing process, which provides sufficient context for answering queries. Additionally, the original text is often highly coupled with less-relevant information, which can introduce noise when used to assist in answering questions.

\subsection{Case Study (RQ3)}
% To facilitate a clear comparison between baseline methods and our \model, we present a specific case in Table~\ref{tab:case}, which includes responses to a machine learning question from both the competitive baseline, GraphRAG, and our \model\ framework. In this instance, \model\ outperforms in all evaluation dimensions assessed by the LLM judge, including comprehensiveness, diversity, empowerment, and overall quality. We make the following observations:
To provide a clear comparison between baseline methods and our \model, we present specific case examples in Table~\ref{tab:case}, which includes responses to a machine learning question from both the competitive baseline, GraphRAG, and our \model\ framework. In this instance, \model\ outperforms in all evaluation dimensions assessed by the LLM judge, including comprehensiveness, diversity, empowerment, and overall quality. Our key observations are as follows:

\textbf{i) Comprehensiveness}. Notably, \model\ covers a broader range of machine learning metrics, showcasing its comprehensiveness and ability to effectively discover relevant information. This highlights the strength of our graph-based indexing paradigm, which excels in precise entity and relation extraction as well as LLM profiling.
% Notably, \model\ discusses a wider range of machine learning metrics, demonstrating its comprehensiveness and its ability to effectively interconnect all pertinent information. This showcases the strength of our graph-based indexing paradigm, which excels in precise entity and relation extraction and LLM profiling.
\textbf{ii) Both Diversity and Empowerment}.
% Furthermore, \model\ not only provides more diverse information but also more empowering content. This success can be attributed to \model's dual-level retrieval paradigm, which conducts in-depth explorations focusing on related entities through low-level retrieval to enhance empowerment and broad explorations via high-level retrieval to improve answer diversity, capturing a comprehensive global perspective of the knowledge domain.
Furthermore, \model\ not only offers a more diverse array of information but also delivers more empowering content. This success is due to \model's hierarchical retrieval paradigm, which combines in-depth explorations of related entities through low-level retrieval to enhance empowerment with broader explorations via high-level retrieval to improve answer diversity. Together, these approaches capture a comprehensive global perspective of the knowledge domain, contributing to better RAG performance.

\begin{table}[t]
\centering
% \caption{A representative case comparing GraphRAG to our \model\ method.}
\caption{Case Study: Comparison Between \model\ and the Baseline Method GraphRAG.}
\vspace{-0.08in}
\label{tab:case}
\resizebox{\textwidth}{!}{
\begin{tabular}{|l|p{15cm}|}
\hline
\textbf{Query} & Which methods can normalize the feature values for effective machine learning? \\ \hline
\textbf{GraphRAG} & \textbf{Key Metrics for Evaluating Movie Recommendation Systems}:  
Evaluating the effectiveness of movie recommendation systems is crucial for ensuring user satisfaction and improving the overall performance of these systems. Several metrics are widely recognized for their ability to provide insights into how well a recommendation system performs. Below are the most informative metrics:  
1. \textbf{Precision} [...]  
2. \textbf{Recall} [...]  
3. \textbf{F1 Score} [...]  
4. \textbf{Mean Average Precision (MAP)} [...]  
5. \textbf{Root Mean Squared Error (RMSE)} [...]  
6. \textbf{User Satisfaction Metrics} [...]  
\textbf{Conclusion}: In summary, metrics such as precision, recall, F1 score, MAP, RMSE, and user satisfaction are critical for evaluating the performance of movie recommendation systems. Each metric provides unique insights that can help developers and researchers refine their algorithms and enhance user experiences. By focusing on these metrics, stakeholders can ensure that their recommendation systems effectively meet user preferences and expectations. \\ \hline
\textbf{LightRAG} & \textbf{Key Metrics for Evaluating Movie Recommendation Systems}:  
Evaluating movie recommendation systems involves multiple metrics that assess the effectiveness and accuracy of the recommendations provided to users. Below are some of the most informative metrics commonly used in this context, derived from the data and concepts surrounding recommendation systems.  
1. \textbf{Mean Average Precision at K (MAPK)} [...]  
2. \textbf{Precision and Recall} [...]  
3. \textbf{Root Mean Squared Error (RMSE) and Mean Squared Error (MSE)} [...]  
4. \textbf{Area Under the Curve (AUC)} [...]  
5. \textbf{F-Measure} [...]  
6. \textbf{User Engagement Metrics} [...]  
\textbf{Conclusion}: A comprehensive evaluation of movie recommendation systems employs a range of metrics that capture both the accuracy of the recommendations and user satisfaction. By effectively leveraging these metrics, system designers can optimize recommendation engines to deliver more personalized and relevant experiences. \\ \hline
\textbf{LLM Decision} & \textbf{Comprehensiveness}:  
"Winner": "Answer 2 (LightRAG)", % \newline

"Explanation": "Answer 2 provides a broader array of metrics including MAPK, AUC, and user engagement metrics alongside precision, recall, RMSE, and F-Measure. This gives a fuller view of how recommendation systems can be evaluated. Answer 1, while detailed, is more focused on traditional metrics and lacks the inclusion of engagement metrics which are crucial for understanding user interaction." \newline

\textbf{Diversity}:  
"Winner": "Answer 2 (LightRAG)", %\newline

"Explanation": "Answer 2 not only covers a wide variety of metrics but also includes nuanced explanations of how some metrics interrelate and differ from one another, like the inclusion of both RMSE and MSE, as well as the details behind AUC. In contrast, Answer 1 sticks primarily to standard metrics without much exploration of potential nuances." \newline

\textbf{Empowerment}:  
"Winner": "Answer 2 (LightRAG)", % \newline

"Explanation": "Answer 2 empowers the reader more effectively by detailing how each metric functions and its importance in evaluating recommendation systems. By providing context such as the trade-offs between precision and recall and emphasizing user engagement metrics, it enables readers to make more informed judgments and understand the implications of different metrics. Answer 1 is more straightforward but lacks the depth of insight regarding why these metrics matter." \newline

% \textbf{Directness}:  
% "Winner": "Answer 1 (GraphRAG)",  \newline

% "Explanation": "Answer 1 presents each metric systematically with clear definitions and implications. It directly addresses the question with straightforward definitions and practical implications for each listed metric. While Answer 2 is also direct, its additional elaborations sometimes lead to a slightly more complex presentation, making it less immediately accessible." \newline

\textbf{Overall Winner}:  
"Winner": "Answer 2 (LightRAG)",  %\newline

"Explanation": "While Answer 1 is more direct and systematic, Answer 2 excels in comprehensiveness, diversity, and empowerment. It provides a richer exploration of the topic, including insights into user engagement and nuanced differences between metrics. This depth and breadth make it more informative for readers seeking to thoroughly understand the evaluation of movie recommendation systems." \\ \hline

\end{tabular}
}
\vspace{-0.2in}
\end{table}

\vspace{-0.05in}
\subsection{Model Cost and Adaptability Analysis (RQ4)}
\vspace{-0.05in}

\begin{wrapfigure}{R}{0.45\textwidth}
    \vspace{-0.18in}
    \centering
    \caption{Comparison of Cost in Terms of Tokens and API Calls for GraphRAG and \model\ on the Legal Dataset.}
    \vspace{-0.1in}
    \label{tab:cost}
    \centering
    \scriptsize
    \setlength{\tabcolsep}{0.3mm}
    % \renewcommand{\arraystretch}{1.5} 
    \begin{tabular}{|c|c|c|c|c|}
        \hline
        Phase & \multicolumn{2}{c|}{Retrieval Phase} & \multicolumn{2}{c|}{Incremental Text Update} \\ \hline
        Model & GraphRAG & Ours & GraphRAG & Ours \\ \hline
        \multirow{2}{*}{Tokens} & \multirow{2}{*}{$610 \times 1{,}000$} & \multirow{2}{*}{$< 100$} & $1{,}399 \times 2 \times 5{,}000$ & \multirow{2}{*}{$T_{\text{extract}}$} \\ 
        &&&$+ T_{\text{extract}}$& \\
        \hline
        API & \multirow{2}{*}{$\frac{610 \times 1{,}000}{C_{\text{max}}}$} & \multirow{2}{*}{1} & \multirow{2}{*}{$1{,}399 \times 2 + C_{\text{extract}}$} & \multirow{2}{*}{$C_{\text{extract}}$}\\
        Calls & &&&\\\hline
        % Calls & $610 \times 1{,}000 \div C_{\text{max}}$ & $1$ & $1{,}399 \times 2 + C_{\text{extract}}$ & $C_{\text{extract}}$ \\ \hline
    \end{tabular}
    \vspace{-0.15in}
\end{wrapfigure}

% \begin{table}[t]
%     \caption{Token Cost and API Calls Comparison for GraphRAG and LightRAG on Legal. 
%     $T_{\text{extract}}$ represents the token overhead for entity and relationship extraction. 
%     $C_{\text{max}}$ represents the maximum number of tokens allowed per API call. 
%     $C_{\text{extract}}$ represents the number of API calls for entity and relationship extraction.}
%     \label{tab:cost}
%     \centering
%     \scriptsize
%     \setlength{\tabcolsep}{0.7mm}
%     % \renewcommand{\arraystretch}{1.5} 
%     \begin{tabular}{|c|c|c|c|c|}
%         \hline
%         Phase & \multicolumn{2}{c|}{Retrieval Phase} & \multicolumn{2}{c|}{Incremental Text Update Phase} \\ \hline
%         Model & GraphRAG & LightRAG & GraphRAG & LightRAG \\ \hline
%         Tokens & $610 \times 1{,}000$ & $< 100$ & $1{,}399 \times 2 \times 5{,}000 + T_{\text{extract}}$ & $T_{\text{extract}}$ \\ \hline
%         API Calls & $610 \times 1{,}000 \div C_{\text{max}}$ & $1$ & $1{,}399 \times 2 + C_{\text{extract}}$ & $C_{\text{extract}}$ \\ \hline
%     \end{tabular}
% \end{table}

% We compare the cost of our \model with that of the representative baseline GraphRAG across two key phases. Firstly, we compare the number of tokens and the number of API calls in the retrieval phase. Secondly, we compare the two metrics in the incremental text update process. The results of this evaluation on the legal dataset is shown in Table~\ref{tab:cost}. Here $T_{\text{extract}}$ represents the token overhead for entity and relationship extraction, $C_{\text{max}}$ represents the maximum number of tokens allowed per API call, and $C_{\text{extract}}$ represents the number of API calls for entity and relationship extraction.
We compare the cost of our \model\ with that of the top-performing baseline, GraphRAG, from two key perspectives. First, we examine the number of tokens and API calls during the indexing and retrieval processes. Second, we analyze these metrics in relation to handling data changes in dynamic environments. The results of this evaluation on the legal dataset are presented in Table~\ref{tab:cost}. In this context, $T_{\text{extract}}$ represents the token overhead for entity and relationship extraction, $C_{\text{max}}$ denotes the maximum number of tokens allowed per API call, and $C_{\text{extract}}$ indicates the number of API calls required for extraction.

% In the retrieval phase, GraphRAG generates 1,399 communities, of which 610 level-2 communities are actively used for retrieval in this experiment. Each community report contains an average of 1,000 tokens, resulting in a total token consumption of 610,000 tokens (610 communities × 1,000 tokens per community). Additionally, due to the need to traverse each community individually, hundreds of API calls are required, significantly increasing the retrieval overhead. In contrast, LightRAG optimizes this process by using less than 100 tokens for keyword generation and retrieval, which requires only a single API call for the entire retrieval process. This efficiency is achieved through our keyword-based retrieval mechanism, eliminating the need to process a large amount of information upfront.
In the retrieval phase, GraphRAG generates 1,399 communities, with 610 level-2 communities actively utilized for retrieval in this experiment. Each community report averages 1,000 tokens, resulting in a total token consumption of 610,000 tokens (610 communities $\times$ 1,000 tokens per community). Additionally, GraphRAG's requirement to traverse each community individually leads to hundreds of API calls, significantly increasing retrieval overhead. In contrast, \model\ optimizes this process by using fewer than 100 tokens for keyword generation and retrieval, requiring only a single API call for the entire process. This efficiency is achieved through our retrieval mechanism, which seamlessly integrates graph structures and vectorized representations for information retrieval, thereby eliminating the need to process large volumes of information upfront.

% In the incremental text update phase, while the entity and relationship extraction overhead is similar for both models, the handling of newly added data shows a stark difference. For a newly added dataset of the same size as the legal, GraphRAG must first deconstruct its existing community structure to incorporate new entities and relationships, and then regenerate the community structure. The token cost for generating each community report is approximately 5,000 tokens. Given the presence of 1,399 communities, GraphRAG would require approximately 1,399 × 2 × 5,000 tokens to reconstruct both original and new community reports. In contrast, LightRAG only integrates newly extracted entities and relationships into the existing graph without needing to reconstruct the entire community structure, resulting in significantly lower overhead during incremental updates.
In the incremental data update phase, designed to address changes in dynamic real-world scenarios, both models exhibit similar overhead for entity and relationship extraction. However, GraphRAG shows significant inefficiency in managing newly added data. When a new dataset of the same size as the legal dataset is introduced, GraphRAG must dismantle its existing community structure to incorporate new entities and relationships, followed by complete regeneration. This process incurs a substantial token cost of approximately 5,000 tokens per community report. Given 1,399 communities, GraphRAG would require around 1,399 × 2 × 5,000 tokens to reconstruct both the original and new community reports—an exorbitant expense that underscores its inefficiency. In contrast, \model\ seamlessly integrates newly extracted entities and relationships into the existing graph without the need for full reconstruction. This approach results in significantly lower overhead during incremental updates, demonstrating its superior efficiency and cost-effectiveness.

